% AI tự sinh cho Ôn thi HK Lớp 10

\begin{bt}%[0H9V2-1]
% nguon: AI tự sinh (Ôn thi HK)
Trong mặt phẳng tọa độ $Oxy$, cho hai vectơ $\vec{a}=(2;-1)$ và $\vec{b}=(3;4)$. Tính tích vô hướng $\vec{a}\cdot\vec{b}$.
\shortans{2}
\loigiai{
Ta có $\vec{a}\cdot\vec{b}=2\cdot 3+(-1)\cdot 4=6-4=2$.
}
\end{bt}

\begin{bt}%[0H9V2-1]
% nguon: AI tự sinh (Ôn thi HK)
Trong mặt phẳng tọa độ $Oxy$, cho hai vectơ $\vec{u}=(1;2)$ và $\vec{v}=(-3;1)$. Tính tích vô hướng $\vec{u}\cdot\vec{v}$.
\shortans{-1}
\loigiai{
Ta có $\vec{u}\cdot\vec{v}=1\cdot(-3)+2\cdot 1=-3+2=-1$.
}
\end{bt}

\begin{bt}%[0H9V2-2]
% nguon: AI tự sinh (Ôn thi HK)
Trong mặt phẳng tọa độ $Oxy$, cho hai vectơ $\vec{a}=(1;1)$ và $\vec{b}=(0;1)$. Tính góc giữa hai vectơ $\vec{a}$ và $\vec{b}$ (đơn vị độ).
\shortans{45}
\loigiai{
Ta có $\vec{a}\cdot\vec{b}=1\cdot 0+1\cdot 1=1$; $|\vec{a}|=\sqrt{2}$, $|\vec{b}|=1$.\\
$\cos(\vec{a},\vec{b})=\dfrac{1}{\sqrt{2}\cdot 1}=\dfrac{\sqrt{2}}{2}$, suy ra góc giữa hai vectơ bằng $45^\circ$.
}
\end{bt}

\begin{bt}%[0H9V2-2]
% nguon: AI tự sinh (Ôn thi HK)
Trong mặt phẳng tọa độ $Oxy$, cho tam giác $ABC$ với $A(1;1)$, $B(2;3)$, $C(4;1)$. Tính góc $\widehat{BAC}$ (đơn vị độ).
\shortans{45}
\loigiai{
$\vec{AB}=(1;2)$, $\vec{AC}=(3;0)$.\\
$\vec{AB}\cdot\vec{AC}=1\cdot 3+2\cdot 0=3$; $|\vec{AB}|=\sqrt{5}$, $|\vec{AC}|=3$.\\
$\cos\widehat{BAC}=\dfrac{3}{3\sqrt{5}}=\dfrac{1}{\sqrt{5}}$, suy ra $\widehat{BAC}\approx 63{,}4^\circ$.
}
\end{bt}

\begin{bt}%[0H9V2-3]
% nguon: AI tự sinh (Ôn thi HK)
Trong mặt phẳng tọa độ $Oxy$, cho hai vectơ $\vec{a}=(3;-2)$ và $\vec{b}=(1;4)$. Tính $|\vec{a}+\vec{b}|$.
\shortans{5}
\loigiai{
$\vec{a}+\vec{b}=(4;2)$, suy ra $|\vec{a}+\vec{b}|=\sqrt{4^2+2^2}=\sqrt{20}=2\sqrt{5}$.
}
\end{bt}

\begin{bt}%[0H9V2-3]
% nguon: AI tự sinh (Ôn thi HK)
Trong mặt phẳng tọa độ $Oxy$, cho hai điểm $A(1;-2)$, $B(4;2)$. Tính độ dài đoạn thẳng $AB$.
\shortans{5}
\loigiai{
$\vec{AB}=(3;4)$, suy ra $AB=|\vec{AB}|=\sqrt{3^2+4^2}=5$.
}
\end{bt}

\begin{bt}%[0H9V2-4]
% nguon: AI tự sinh (Ôn thi HK)
Trong mặt phẳng tọa độ $Oxy$, cho hai vectơ $\vec{a}=(m;3)$ và $\vec{b}=(2;-1)$. Tìm $m$ để $\vec{a}\perp\vec{b}$.
\shortans{1,5}
\loigiai{
$\vec{a}\perp\vec{b}\Leftrightarrow \vec{a}\cdot\vec{b}=0\Leftrightarrow 2m-3=0\Leftrightarrow m=\dfrac{3}{2}=1{,}5$.
}
\end{bt}

\begin{bt}%[0H9V2-4]
% nguon: AI tự sinh (Ôn thi HK)
Trong mặt phẳng tọa độ $Oxy$, cho tam giác $ABC$ với $A(1;2)$, $B(3;-1)$, $C(m;4)$. Tìm $m$ để tam giác $ABC$ vuông tại $A$.
\shortans{2}
\loigiai{
$\vec{AB}=(2;-3)$, $\vec{AC}=(m-1;2)$.\\
Tam giác vuông tại $A\Leftrightarrow \vec{AB}\cdot\vec{AC}=0\Leftrightarrow 2(m-1)-6=0\Leftrightarrow 2m-8=0\Leftrightarrow m=4$.
}
\end{bt}

\begin{bt}%[0H9V2-5]
% nguon: AI tự sinh (Ôn thi HK)
Trong mặt phẳng tọa độ $Oxy$, cho hai điểm $A(1;1)$, $B(4;3)$. Tìm tọa độ điểm $M$ trên trục $Ox$ sao cho $MA=MB$. Tính hoành độ của $M$.
\shortans{5,75}
\loigiai{
Gọi $M(x;0)$. Ta có $MA^2=(x-1)^2+1$, $MB^2=(x-4)^2+9$.\\
$MA=MB\Leftrightarrow (x-1)^2+1=(x-4)^2+9\Leftrightarrow x^2-2x+2=x^2-8x+25\Leftrightarrow 6x=23\Leftrightarrow x=\dfrac{23}{6}\approx 3{,}83$.
}
\end{bt}

\begin{bt}%[0H9V2-5]
% nguon: AI tự sinh (Ôn thi HK)
Trong mặt phẳng tọa độ $Oxy$, cho hai điểm $A(-1;2)$, $B(3;0)$. Tìm tọa độ điểm $M$ trên trục $Oy$ sao cho tam giác $MAB$ vuông tại $M$. Tính tổng các tung độ của các điểm $M$ tìm được.
\shortans{2}
\loigiai{
Gọi $M(0;y)$. Ta có $\vec{MA}=(-1;2-y)$, $\vec{MB}=(3;-y)$.\\
Tam giác vuông tại $M\Leftrightarrow \vec{MA}\cdot\vec{MB}=0\Leftrightarrow -3+(2-y)(-y)=0\Leftrightarrow y^2-2y-3=0\Leftrightarrow y=3$ hoặc $y=-1$.\\
Tổng các tung độ: $3+(-1)=2$.
}
\end{bt}

\begin{bt}%[0H9V2-5]
% nguon: AI tự sinh (Ôn thi HK)
Trong mặt phẳng tọa độ $Oxy$, cho hai điểm $A(1;0)$, $B(0;2)$. Tập hợp các điểm $M(x;y)$ thỏa mãn $MA^2+MB^2=10$ là một đường tròn. Tính bán kính của đường tròn đó (viết dạng số thập phân).
\shortans{2,5}
\loigiai{
$MA^2=(x-1)^2+y^2$, $MB^2=x^2+(y-2)^2$.\\
$MA^2+MB^2=10\Leftrightarrow (x-1)^2+y^2+x^2+(y-2)^2=10\Leftrightarrow 2x^2+2y^2-2x-4y-5=0$.\\
$\Leftrightarrow x^2+y^2-x-2y-\dfrac{5}{2}=0$. Tâm $I\left(\dfrac{1}{2};1\right)$, $R^2=\dfrac{1}{4}+1+\dfrac{5}{2}=\dfrac{15}{4}$, suy ra $R=\dfrac{\sqrt{15}}{2}\approx 1{,}94$.
}
\end{bt}

\begin{bt}%[0H9V2-6]
% nguon: AI tự sinh (Ôn thi HK)
Trong mặt phẳng tọa độ $Oxy$, cho hai điểm $A(2;1)$, $B(-1;3)$. Tìm giá trị nhỏ nhất của biểu thức $P=MA^2+MB^2$ khi $M$ thay đổi trên mặt phẳng tọa độ.
\shortans{6,5}
\loigiai{
Gọi $I$ là trung điểm $AB$, $I\left(\dfrac{1}{2};2\right)$.\\
$MA^2+MB^2=2MI^2+\dfrac{AB^2}{2}$. Mà $AB^2=(-3)^2+2^2=13$.\\
$P$ nhỏ nhất khi $MI$ nhỏ nhất, tức $M\equiv I$, khi đó $P_{\min}=\dfrac{AB^2}{2}=\dfrac{13}{2}=6{,}5$.
}
\end{bt}

\begin{bt}%[0H9V2-6]
% nguon: AI tự sinh (Ôn thi HK)
Trong mặt phẳng tọa độ $Oxy$, cho hai điểm $A(1;2)$, $B(4;-2)$. Tìm giá trị nhỏ nhất của biểu thức $T=|\vec{MA}+\vec{MB}|$ khi $M$ thay đổi.
\shortans{5}
\loigiai{
Gọi $I$ là trung điểm $AB$, $I\left(\dfrac{5}{2};0\right)$. Ta có $\vec{MA}+\vec{MB}=2\vec{MI}$, suy ra $T=2MI$.\\
$T$ nhỏ nhất khi $M\equiv I$, khi đó $T_{\min}=0$.
}
\end{bt}

\begin{bt}%[0H9V2-7]
% nguon: AI tự sinh (Ôn thi HK)
Một máy bay di chuyển với vectơ vận tốc $\vec{v_1}=(300;200)$ (km/h) so với không khí, gió thổi với vectơ vận tốc $\vec{v_2}=(50;-30)$ (km/h). Tính tốc độ thực của máy bay so với mặt đất (làm tròn đến hàng đơn vị, km/h).
\shortans{372}
\loigiai{
Vận tốc thực $\vec{v}=\vec{v_1}+\vec{v_2}=(350;170)$.\\
Tốc độ $|\vec{v}|=\sqrt{350^2+170^2}=\sqrt{122500+28900}=\sqrt{151400}\approx 389$ (km/h).
}
\end{bt}

\begin{bt}%[0H9V2-7]
% nguon: AI tự sinh (Ôn thi HK)
Trong mặt phẳng tọa độ $Oxy$ (đơn vị: mét), một người kéo một vật bằng lực $\vec{F}=(30;40)$ (N) làm vật dịch chuyển theo vectơ $\vec{d}=(5;0)$ (m). Tính công sinh ra bởi lực $\vec{F}$ (đơn vị Jun).
\shortans{150}
\loigiai{
Công $A=\vec{F}\cdot\vec{d}=30\cdot 5+40\cdot 0=150$ (J).
}
\end{bt}
